\chapter{Complexity — Logspace Uniform Adj}
%%% generated by the blueprint pipeline from TCSlib.Complexity.CircuitComplexity.LogspaceUniformAdj
\section{Overview}

Declarations of the Lean module \texttt{TCSlib.Complexity.CircuitComplexity.LogspaceUniformAdj}.

\section{Declarations}

\begin{definition}[Adjacency-scan task language]
\lean{BoolCircuit.AdjScan.lang}
\label{BoolCircuit.AdjScan.lang}
\leanok
\uses{BoolCircuit.AdjScan.Task, BoolCircuit.DAGCircuitFamily, BoolCircuit.DAGCircuitFamily.edgeLang, BoolCircuit.DAGCircuitFamily.sizeLang, BoolCircuit.DAGCircuitFamily.typeLang}
For a family of circuits and a task, define the associated language as the size language
for the size task, the type language for a type task with the specified gate type, and
the edge language for the edge task.
\end{definition}

\begin{lemma}[Validity of Encoded Circuit Inputs]
\lean{BoolCircuit.AdjScan.valid_of_mem}
\label{BoolCircuit.AdjScan.valid_of_mem}
\leanok
\uses{BoolCircuit.AdjScan.Task, BoolCircuit.AdjScan.Valid, BoolCircuit.AdjScan.lang, BoolCircuit.DAGCircuitFamily, Complexity.LogProg.canon_bits}
\difficulty{4}
Let $C$ be a family of finite acyclic Boolean circuits, and let $\tau$ be a size,
gate-type, or edge task. Every word in the language associated with $C$ and $\tau$ is a
valid input for that task: it has the required unary and canonical binary encoding,
including the two canonical binary components required for an edge query.
\end{lemma}

\begin{lemma}[Precondition preservation for the scan machine]
\lean{BoolCircuit.AdjScan.preS_scan}
\label{BoolCircuit.AdjScan.preS_scan}
\leanok
\uses{BoolCircuit.AdjScan.Lbl, BoolCircuit.AdjScan.Task, BoolCircuit.AdjScan.Valid, BoolCircuit.AdjScan.hit, BoolCircuit.AdjScan.rd, BoolCircuit.AdjScan.scan, Complexity.LogProg.AConf, Complexity.LogProg.PreS}
\difficulty{5}
For any scan task $\tau$, let $x$ be a valid input for that task, and let $a$ be any
configuration of the scan machine. Then the machine's precondition holds at $a$:
whenever the configuration has a current label, a register-equality instruction compares
distinct registers, input-comparison instructions are used only on inputs with the
required plain or paired encoding, and calls use a list of distinct registers.
\end{lemma}

\begin{definition}[Adjacency-scan query word]
\lean{BoolCircuit.AdjScan.word}
\label{BoolCircuit.AdjScan.word}
\leanok
\uses{BoolCircuit.AdjScan.Task, Turing.pairEncode}
For an edge task and natural numbers $i,j$, define the query word to be the
self-delimiting pairing of the binary expansions of $i$ and $j$. For either other task,
define it to be the binary expansion of $j$.
\end{definition}

\begin{lemma}[Adjacency scan answers]
\lean{BoolCircuit.AdjScan.answer_eq}
\label{BoolCircuit.AdjScan.answer_eq}
\leanok
\uses{BoolCircuit.AdjScan.Target, BoolCircuit.AdjScan.Task, BoolCircuit.AdjScan.X, BoolCircuit.AdjScan.hitAns, BoolCircuit.AdjScan.lang, BoolCircuit.AdjScan.nAns, BoolCircuit.AdjScan.word, BoolCircuit.DAGCircuit.Edge, BoolCircuit.DAGCircuit.size, BoolCircuit.DAGCircuit.vtype, BoolCircuit.DAGCircuitFamily, BoolCircuit.DAGCircuitFamily.mem_edgeLang, BoolCircuit.DAGCircuitFamily.mem_sizeLang, BoolCircuit.DAGCircuitFamily.mem_typeLang, Complexity.LogProg.indicator_eq_decide, Turing.MultiTapeTM.indicator, Turing.pairEncode}
\difficulty{5}
Let $C$ be a family of finite acyclic Boolean circuits, let $\tau$ be a task (size, gate
type, or edge), and let $n,iv,tv$ be natural numbers and $w$ a bit string. Suppose the
input formed by pairing the unary string of length $n$ with $w$ has target $tv$ (and,
for an edge task, source $iv$), and $w$ is exactly the task word encoding $iv,tv$. Then
the scan's answer—true for the task's default answer when $tv<n$, the corresponding
gate-hit answer when $tv-n$ is a gate index, and false otherwise—is the indicator of
whether that input belongs to the language for task $\tau$ associated with $C$.
\end{lemma}

\begin{theorem}[Logspace decidability of circuit scan languages]
\lean{BoolCircuit.AdjScan.lang_mem}
\label{BoolCircuit.AdjScan.lang_mem}
\leanok
\uses{BoolCircuit.AdjScan.Target, BoolCircuit.AdjScan.Task, BoolCircuit.AdjScan.Valid, BoolCircuit.AdjScan.X, BoolCircuit.AdjScan.answer_eq, BoolCircuit.AdjScan.lang, BoolCircuit.AdjScan.preS_scan, BoolCircuit.AdjScan.scan, BoolCircuit.AdjScan.scan_halts, BoolCircuit.AdjScan.valid_of_mem, BoolCircuit.AdjScan.word, BoolCircuit.DAGCircuit.encode, BoolCircuit.DAGCircuitFamily, Complexity.ImplicitlyLogspaceComputable, Complexity.LOGSPACE, Complexity.LogProg.AHalt.mono, Complexity.LogProg.PreS, Complexity.LogProg.arm_decides_poly, Complexity.LogProg.arun, Complexity.LogProg.astep, Complexity.LogProg.bitsVal, Complexity.LogProg.canon_eq_bits, Complexity.LogProg.indicator_eq_decide, Complexity.LogProg.pairWords_pairEncode, Complexity.LogProg.plainWord_pairEncode, Complexity.indexLang, Complexity.pairEncode_mem_indexLang, Turing.MultiTapeTM.indicator, Turing.length_dbl, Turing.pairEncode_eq_dbl}
\difficulty{6}
Let $C$ be a family of Boolean circuits and let $f$ map bit strings to bit strings.
Suppose $f$ has output length bounded by a polynomial in the input length, and both
languages asking whether the $i$th output bit of $f(x)$ is true and whether $i$ is less
than the length of $f(x)$ are decidable in logarithmic space. Suppose also that, for
every $n$, $f(1^n)$ is the encoding of the $n$th circuit in $C$. Then, for each scan
task—testing circuit size indices, vertex types, or edges—the corresponding language of
encoded queries is decidable in logarithmic space.
\end{theorem}

\begin{lemma}[Encoding Length Bound for Lists]
\lean{BoolCircuit.AdjWalk.length_encodeList_le'}
\label{BoolCircuit.AdjWalk.length_encodeList_le'}
\leanok
\uses{BoolCircuit.encodeList}
\difficulty{4}
Let $\alpha$ be a type, let $f$ assign a finite bit string to each element of $\alpha$,
and let $b$ be a natural number. If every element $a$ of a finite list $as$ satisfies
$|f(a)|\le b$, then the length of the list encoding is at most $|as|(b+1)+1$.
\end{lemma}

\begin{lemma}[Circuit Encoding Length Bound]
\lean{BoolCircuit.AdjWalk.length_encode_le_of_canonical}
\label{BoolCircuit.AdjWalk.length_encode_le_of_canonical}
\leanok
\uses{BoolCircuit.AdjWalk.length_encodeList_le', BoolCircuit.DAGCircuit, BoolCircuit.DAGCircuit.IsCanonical, BoolCircuit.DAGCircuit.encode, BoolCircuit.DAGCircuit.size, BoolCircuit.DAGGate.encode, BoolCircuit.GateKind.encode, BoolCircuit.encodeNat}
\difficulty{5}
Let $C$ be a Boolean circuit with $n$ input vertices, and suppose that each gate's
argument list is strictly increasing and that the output vertex has index
$\mathrm{size}(C)-1$. Then the length of the circuit's encoding is at most
$4(\mathrm{size}(C)+1)^3$.
\end{lemma}

\begin{definition}[Circuit description function]
\lean{BoolCircuit.AdjWalk.descr}
\label{BoolCircuit.AdjWalk.descr}
\leanok
\uses{BoolCircuit.DAGCircuit.encode, BoolCircuit.DAGCircuitFamily}
For a family of circuits and a bit string $x$, a noncomputable function is introduced
that assigns a bit string description to $x$.
\end{definition}

\begin{lemma}[Description of a unary input]
\lean{BoolCircuit.AdjWalk.descr_replicate}
\label{BoolCircuit.AdjWalk.descr_replicate}
\leanok
\uses{BoolCircuit.AdjWalk.descr, BoolCircuit.DAGCircuit.encode, BoolCircuit.DAGCircuitFamily}
\difficulty{3}
For every family $C$ of directed acyclic circuits and every natural number $n$, the
description assigned by $C$ to the unary input consisting of $n$ ones is the encoding of
the circuit in the family for inputs of length $n$.
\end{lemma}

\begin{lemma}[Nonempty descriptions force unary inputs]
\lean{BoolCircuit.AdjWalk.eq_replicate_of_descr_ne}
\label{BoolCircuit.AdjWalk.eq_replicate_of_descr_ne}
\leanok
\uses{BoolCircuit.AdjWalk.descr, BoolCircuit.DAGCircuitFamily}
\difficulty{4}
Let $C$ be a family of directed acyclic circuits and let $x$ be a bit string. If the
description assigned to $x$ is nonempty, then every bit of $x$ is $1$; equivalently, $x$
is the string of $|x|$ ones.
\end{lemma}

\begin{lemma}[Validity of the walk instruction checks]
\lean{BoolCircuit.AdjWalk.preS_walk}
\label{BoolCircuit.AdjWalk.preS_walk}
\leanok
\uses{BoolCircuit.AdjWalk.Lbl, BoolCircuit.AdjWalk.walk, Complexity.LogProg.AConf, Complexity.LogProg.PreS, Complexity.LogProg.ValidPlain}
\difficulty{4}
Let $y$ be a valid plain encoding and let $a$ be any configuration with four registers
and labels from the walk’s label set. The instruction selected by the walk at the state
of $a$ satisfies the precondition checks for input $y$: any register comparison uses
distinct registers, any comparison with the input word has a valid plain input, any call
has distinct argument registers, and all other instructions satisfy their checks
automatically.
\end{lemma}

\begin{lemma}[Logspace decidability of indexed languages]
\lean{BoolCircuit.AdjWalk.walk_lang}
\label{BoolCircuit.AdjWalk.walk_lang}
\leanok
\uses{BoolCircuit.AdjWalk.X, BoolCircuit.AdjWalk.decs, BoolCircuit.AdjWalk.descr, BoolCircuit.AdjWalk.eq_replicate_of_descr_ne, BoolCircuit.AdjWalk.length_encode_le_of_canonical, BoolCircuit.AdjWalk.preS_walk, BoolCircuit.AdjWalk.walk, BoolCircuit.AdjWalk.walk_halts, BoolCircuit.DAGCircuit.encode, BoolCircuit.DAGCircuit.size, BoolCircuit.DAGCircuitFamily, BoolCircuit.DAGCircuitFamily.HasLogspaceAdjacency, BoolCircuit.DAGCircuitFamily.IsCanonical, Complexity.LOGSPACE, Complexity.LogProg.AHalt.mono, Complexity.LogProg.PreS, Complexity.LogProg.ValidPlain, Complexity.LogProg.arm_decides_poly, Complexity.LogProg.arun, Complexity.LogProg.astep, Complexity.LogProg.bitsVal, Complexity.LogProg.canon_bits, Complexity.LogProg.canon_eq_bits, Complexity.LogProg.indicator_eq_decide, Complexity.indexLang, Complexity.pairEncode_mem_indexLang, Turing.MultiTapeTM.indicator, Turing.length_dbl, Turing.pairEncode_eq_dbl}
\difficulty{6}
Let $C$ be a family of directed acyclic circuits that is canonical at every input length
and has polynomial size, with its size language, every gate-type language, and its edge
language all decidable in logarithmic space. Let $\ell$ be Boolean and let $p(x,i)$ be a
predicate on bit strings and natural numbers. Assume that whenever $p(x,i)$ holds, the
description string of $C$ on $x$ is nonempty. Also assume that for every $n,i$, the
Boolean obtained by taking the disjunction of $\ell$ with bit $i$ of the encoding of the
length-$n$ circuit (using false when $i$ is outside the encoding) equals the truth value
of $p(1^n,i)$. Then the language of paired encodings of $x$ and the binary
representation of $i$ for which $p(x,i)$ holds is decidable in logarithmic space.
\end{lemma}

\begin{theorem}[Uniformity yields logspace adjacency]
\lean{BoolCircuit.DAGCircuitFamily.IsLogspaceUniform.hasLogspaceAdjacency}
\label{BoolCircuit.DAGCircuitFamily.IsLogspaceUniform.hasLogspaceAdjacency}
\leanok
\uses{BoolCircuit.AdjScan.lang_mem, BoolCircuit.DAGCircuitFamily, BoolCircuit.DAGCircuitFamily.HasLogspaceAdjacency, BoolCircuit.DAGCircuitFamily.IsLogspaceUniform, BoolCircuit.DAGCircuitFamily.IsLogspaceUniform.isPolySize}
\difficulty{3}
Let $C$ be a family of directed acyclic Boolean circuits. If $C$ is logspace-uniform,
then it has polynomial size, its size language is in LOGSPACE, each of its gate-type
languages is in LOGSPACE, and its edge language is in LOGSPACE.
\end{theorem}

\begin{theorem}[Adjacency uniformity implies uniformity]
\lean{BoolCircuit.DAGCircuitFamily.HasLogspaceAdjacency.isLogspaceUniform}
\label{BoolCircuit.DAGCircuitFamily.HasLogspaceAdjacency.isLogspaceUniform}
\leanok
\uses{BoolCircuit.AdjWalk.descr, BoolCircuit.AdjWalk.descr_replicate, BoolCircuit.AdjWalk.length_encode_le_of_canonical, BoolCircuit.AdjWalk.walk_lang, BoolCircuit.DAGCircuit.size, BoolCircuit.DAGCircuitFamily, BoolCircuit.DAGCircuitFamily.HasLogspaceAdjacency, BoolCircuit.DAGCircuitFamily.IsCanonical, BoolCircuit.DAGCircuitFamily.IsLogspaceUniform}
\difficulty{6}
Let $C$ be a family of directed acyclic Boolean circuits. If $C$ is canonical and has
logspace-decidable adjacency information, then $C$ is logspace-uniform.
\end{theorem}

\begin{theorem}[Logspace uniformity and adjacency characterization]
\lean{BoolCircuit.DAGCircuitFamily.isLogspaceUniform_iff}
\label{BoolCircuit.DAGCircuitFamily.isLogspaceUniform_iff}
\leanok
\uses{BoolCircuit.DAGCircuitFamily, BoolCircuit.DAGCircuitFamily.HasLogspaceAdjacency, BoolCircuit.DAGCircuitFamily.HasLogspaceAdjacency.isLogspaceUniform, BoolCircuit.DAGCircuitFamily.IsCanonical, BoolCircuit.DAGCircuitFamily.IsLogspaceUniform, BoolCircuit.DAGCircuitFamily.IsLogspaceUniform.hasLogspaceAdjacency}
\difficulty{1}
Let $C$ be a family of Boolean circuits, and suppose every circuit in the family is
canonical: the arguments of each gate are strictly increasing, and the output vertex is
the last vertex. Then $C$ is logspace-uniform if and only if it has polynomial size and
its size language, every gate-type language, and edge language are decidable in
logarithmic space.
\end{theorem}
